\documentclass[a4paper,12pt]{article}
\usepackage[utf8]{inputenc}
\usepackage{parskip}
\title{Material}
\date{}
\begin{document}
\maketitle
\section*{Att skriva ut}

I det här dokumentet beskrivs hur mycket av vad som ska skrivas ut. Eftersom antalet grupper varierar från klass till klass kan pappersmängden också behöva justeras. Det kan också vara bra att skriva ut något extra papper ifall något går sönder eller tappas bort.

\textit{Om en kommentar börjar med *, innebär det att eleverna kommer skriva på pappret och att det behövs skrivas ut nya för varje klass.}

\textit{En mapp} \texttt{mapp/} \textit{följt av en stjärna} \texttt{*}\textit{, har betydelsen: alla filer som ligger under} \texttt{mapp/}\textit{.}

\subsection*{Inledning}

\begin{tabular}{|l|l|p{6cm}|}
\hline
\textbf{Filnamn} & \textbf{Antal} & \textbf{Kommentar} \\
\hline
ANSTÄLLNINGSKONTRAKT.pdf  & 3-6st & *En per grupp \\
\hline
storyhäfte\_inledning.pdf & 3-6st & En per grupp \\
\hline
Attackprofiler/*          & 1st   & Varje grupp får sin egen attackprofil \\
\hline
\end{tabular}

\subsection*{Akt 1}

\begin{tabular}{|l|l|p{6cm}|}
\hline
\textbf{Filnamn} & \textbf{Antal} & \textbf{Kommentar} \\
\hline
MyndighetsKlassificering.pdf & 3-6st & *En per grupp \\
\hline
storyhäfte\_akt\_1.pdf       & 3-6st & En per grupp \\
\hline
\end{tabular}

\newpage

\subsection*{Akt 2}

Om möjligt kan övningar som bara har ett blad för övningen och ett blad för uppgiftsbeskrivningen kunna vara på ett tvåsidigt papper.

\begin{tabular}{|l|l|p{6cm}|}
\hline
\textbf{Filnamn} & \textbf{Antal} & \textbf{Kommentar} \\
\hline
storyhäfte\_akt\_2.pdf                  & 3-6st & En per grupp \\
\hline
DDoS uppgiftsbeskrivning.pdf            & 2st   & Tillhör DDOS \\
\hline
DDoS serverlogg.pdf                     & 2st   & *Tillhör DDOS \\
\hline
Deepfake uppgiftsbeskrivning.pdf        & 2st   & Tillhör DEEPFAKE \\
\hline
Deepfake uppgift.pdf                    & 2st   & *Tillhör DEEPFAKE \\
\hline
Domain hijack uppgiftsbeskrivning.pdf   & 2st   & Tillhör DOMAIN HIJACK \\
\hline
Momoni gov a.pdf                        & 2st   & Tillhör DOMAIN HIJACK \\
\hline
Momoni gov b.pdf                        & 2st   & Tillhör DOMAIN HIJACK \\
\hline
Momoni gov c riktig.pdf                 & 2st   & Tillhör DOMAIN HIJACK \\
\hline
Momoni gov d.pdf                        & 2st   & Tillhör DOMAIN HIJACK \\
\hline
Phishing uppgiftsbeskrivning.pdf        & 2st   & Tillhör PHISHING \\
\hline
Phishing mejl.pdf                       & 2st   & *Tillhör PHISHING \\
\hline
Ransomware uppgiftsbeskrivning.pdf      & 2st   & Tillhör RANSOMWARE \\
\hline
Ransomware k1ngp1n.pdf                  & 2st   & *Tillhör RANSOMWARE \\
\hline
Worm uppgiftsbeskrivning.pdf            & 2st   & Tillhör WORM \\
\hline
Worm Attack.pdf                         & 2st   & *Tillhör WORM \\
\hline
\end{tabular}

\subsection*{Akt 3}

\begin{tabular}{|l|l|p{6cm}|}
\hline
\textbf{Filnamn} & \textbf{Antal} & \textbf{Kommentar} \\
\hline
storyhäfte\_akt\_3.pdf & 3-6st & En per grupp \\
\hline
\end{tabular}

\end{document}
